\section{Conclusiones}
    En este miniproyecto reforcé desde la perspectiva práctica y
    tácita como las transformaciones vistas en clase pueden alterar 
    o cambiar cómo se ve una imagen de forma drástica (tanto para 
    bien como para mal), y al implementar manualmente estas transformaciones 
    me puedo dar cuenta que mejorar una imagen no es complicado 
    algorítmicamente, sino encontrar los parámetros más convenientes para 
    cada caso y objetivo. Volviéndose más un arte que una metodología 
    rígida o explícita, claro que se puede comenzar por emplear los 
    histogramas y la escala de grises de una imagen, pero solamente 
    dan un pauta sobre cuáles son las técnicas más adecuadas. Dicho de 
    otra forma, el embellecimiento de una imagen recae más en la 
    experiencia que tiene el editor (humano) y que, para este caso, 
    adquirí un poco de sensibilidad o intuición sobre cómo los 
    parámetros y algoritmos alteran los elementos visuales (tonalidades, 
    contraste, brillo) en una composición, haciendo que el proyecto 
    sea más ameno para su desarrollo y análisis.\\

    El interiorizar estas ideas junto con buenas prácticas de programación 
    me permitió entender de mejor manera cómo los sistemas de imágenes 
    actúan al momento de postprocesar y que detrás de estos mecanismos 
    existen, mayormente, humanos que ajustan los parámetros de las técnicas 
    seleccionadas para mejorar la calidad final y reducir el ruido 
    visual con el fin de resaltar los elementos u objetos de interés 
    para algún proceso posterior, siendo estos procesos los que otorgan 
    un valor a cada imagen con base a qué tan bien fue procesada y tratada. 
    Este aspecto es lo que principal me deja este proyecto, es decir, 
    el valor escondido que adquiere una imagen cuando esta se le aplica 
    las técnicas de mejoramiento adecuadas, donde el valor puede ser 
    información importante para una empresa o facilitar la ejecución 
    o funcionamiento de otros algoritmos más sofisticados.